% Going back to the original firmware on the Samsung YP-S3. Summary of
% "3. Go back to the original firmware" in YPS3-PORT.md.
\label{ref:yps3_uninstall}
Going back to the original firmware works like the installation: the
script \fname{yps3-restore.sh} in \fname{utils/s5l8700\_dfu/} copies the original
firmware image (by default \fname{S3-original.ROM}) to the music drive as
\fname{S3.ROM}, and the \playerman{} bootloader writes it at the next
power-on. \fname{YPS3-PORT.md} in the source tree has the exact commands.

\begin{description}
  \item[Through Rockbox's USB drive (the normal way).]\leavevmode
    \begin{enumerate}
      \item Start Rockbox and connect the \dap{} to the computer.
      \item Run \fname{yps3-restore.sh MUSIC\_DRIVE}.
      \item Eject the drive, unplug the cable, switch the \dap{} off and on.
        The original firmware starts; its settings may be reset. You can
        delete the \fname{.rockbox} directory from the drive afterwards.
    \end{enumerate}
  \item[Through MTP, when the original firmware starts but Rockbox doesn't.]\leavevmode
    \begin{enumerate}
      \item Start the original firmware (keep the \ButtonPower{} slider pushed
        through the start animation, see \reference{ref:Dualboot}), connect
        the \dap{} and open it once in the file manager so that it gets
        mounted.
      \item \fname{gio mount -l} shows the \daps{} \fname{mtp://}
        address. Its storage folder \fname{S3} is the top folder of the music drive.
      \item Run \fname{yps3-restore.sh -{}-mtp mtp://<device>/S3/}.
      \item Disconnect in the file manager, unplug the cable, switch the
        \dap{} off and on. A short upgrade screen appears, then the original
        firmware starts; \fname{S3.ROM} is gone afterwards. If not, use DFU
        mode (see \reference{ref:yps3_rescue}).
    \end{enumerate}
\end{description}
